%%%%%%%%%%%%%%%%%%%%%%%%%%%%%%%%%%%%%%%%%%%%%%%%%%%%%%%%%%%%%%
%%%%%%%%%%%%%%%%%% MA-0806 Análisis funcional %%%%%%%%%%%%%%%%
%%%%%%%%%%%%%%%%%%%%%%%%%%%%%%%%%%%%%%%%%%%%%%%%%%%%%%%%%%%%%%
% Propuesta de programa. Fuentes resumidas y adaptadas:
% - NYU Courant, MATH-GA 2550 Functional Analysis I
% - UC Berkeley, MATH 202B (bloque de análisis funcional) y MATH 206
%   (temas de segundo curso, aquí como ampliación)
% - Georgia Tech, MATH 6338 Real Analysis II y MATH 7338 Functional Analysis
%   (este último, como ampliación)
\newpage
\subsection*{MA-0806 Análisis funcional}
\addcontentsline{toc}{subsection}{MA-0806 Análisis funcional}

\hypertarget{MA0806}{}
\begin{refsection}
\begin{table}[H]
\centering{
\begin{tabular}{|ll|}
\hline
\textbf{Año:} IV  & \textbf{Requisitos:} MA-0705 (o MA-0625) Análisis Real I \\
\textbf{Ciclo:} Optativa  & \textbf{Correquisitos:} Ninguno \\
\textbf{Tipo de curso:} Teórico & \textbf{Horas presenciales por semana:} 5 \\
\textbf{Créditos:} 4 & \textbf{Horas de trabajo independiente por semana:} 7 \\ \hline
\end{tabular}
}
\end{table}

\subsubsection*{Descripción del curso}

Este curso optativo introduce el análisis funcional como el estudio de espacios de dimensión infinita ---principalmente espacios de Banach y de Hilbert--- y de los operadores lineales acotados entre ellos. Se ubica en el cuarto año, una vez que la persona estudiante dispone de la medida e integración de Lebesgue y de una primera aproximación a los espacios $L^p$ (MA-0705 Análisis Real~I). Marca el paso del análisis en $\R^n$ al análisis en espacios abstractos, y constituye una base para estudios posteriores en ecuaciones diferenciales, análisis armónico, teoría espectral y análisis aplicado.

El programa se articula a partir de referencias curriculares contemporáneas e incluye temas como espacios $L^p$, Banach y Hilbert; operadores acotados; teoremas de Hahn--Banach, acotación uniforme, aplicación abierta y gráfica cerrada, dualidad, topologías débiles, teorema de Alaoglu, operadores compactos y teorema espectral.

El énfasis es teórico: se requieren definiciones precisas, enunciados completos y demostraciones de los teoremas fundamentales. Los ejemplos en $L^p$, $C(K)$, $\ell^p$ y espacios de Hilbert clásicos se usan para concretar la teoría, no para sustituirla. Se espera que la persona estudiante formule argumentos rigurosos, construya ejemplos y contraejemplos, y reconozca el alcance de cada hipótesis.

\subsubsection*{Objetivos}

\textbf{General:} Que la persona estudiante sea capaz de resolver problemas y desarrollar teoría a partir de los fundamentos del análisis funcional en espacios de Banach y de Hilbert.

\textbf{Específicos:} Al finalizar el curso se espera que la persona estudiante sea capaz de:

\begin{enumerate}
\item Reconocer y comparar las estructuras de espacio normado, de Banach y de Hilbert, y verificar completitud, separabilidad y reflexividad en ejemplos clásicos, en particular en $L^p$ y $\ell^p$.
\item Emplear las desigualdades de Hölder, Minkowski y Bessel, y el teorema de la proyección, para resolver problemas en espacios de Hilbert y justificar la existencia de bases ortonormales.
\item Analizar operadores lineales acotados, calcular normas de operadores, y relacionar un operador con su adjunto mediante el teorema de representación de Riesz en Hilbert.
\item Demostrar el teorema de Hahn--Banach y aplicarlo a la separación de convexos, a la dualidad de $L^p$ y a la extensión de funcionales.
\item Deducir, a partir del teorema de Baire, los principios de acotación uniforme, de la aplicación abierta y de la gráfica cerrada, e identificar hipótesis indispensables mediante contraejemplos.
\item Distinguir las topologías fuerte, débil y débil-$*$, y aplicar el teorema de Banach--Alaoglu a problemas de existencia y compactidad.
\item Caracterizar operadores compactos y demostrar el teorema espectral para operadores compactos autoadjuntos en Hilbert, con una aplicación a ecuaciones integrales.
\item Consultar literatura especializada y comunicar argumentos de análisis funcional con rigor, tanto de forma oral como escrita.
\end{enumerate}

\subsubsection*{Contenidos}

\begin{enumerate}
\item \textbf{Espacios de Banach y espacios $L^p$ (2 semanas):}
\begin{enumerate}
    \item Espacios normados y de Banach. Equivalencia de normas en dimensión finita. Completitud y series absolutamente convergentes.
    \item Ejemplos: $C(K)$, $C_0$, $\ell^p$, $L^p(\mu)$ y operadores de multiplicación o de desplazamiento.
    \item Desigualdades de Hölder, Minkowski y Jensen. Completitud de $L^p$ y de $\ell^p$.
    \item Densidad de funciones continuas o simples, separabilidad de $L^p$ para $1\le p<\infty$, y el caso $L^\infty$.
\end{enumerate}

\item \textbf{Espacios de Hilbert (2 semanas):}
\begin{enumerate}
    \item Productos internos, identidad del paralelogramo y caracterización de normas hilbertianas.
    \item Ortogonalidad, teorema de la proyección y teorema de la representación de Riesz.
    \item Sistemas ortonormales, desigualdad de Bessel, identidad de Parseval y bases de Hilbert. El espacio $L^2$ como prototipo.
\end{enumerate}

\item \textbf{Operadores lineales acotados (2 semanas):}
\begin{enumerate}
    \item Operadores acotados y continuidad. Norma de operador y el álgebra $\mathcal{B}(X,Y)$.
    \item Adjunto de un operador en Hilbert. Operadores autoadjuntos, unitarios, normales y de Hilbert--Schmidt (nociones).
    \item Espectro y resolvente de un operador acotado: definiciones y propiedades elementales (espectro no vacío en el caso complejo).
\end{enumerate}

\item \textbf{Teoremas fundamentales (4 semanas):}
\begin{enumerate}
    \item Teorema de Baire y consecuencias para espacios de Banach.
    \item Principio de acotación uniforme (Banach--Steinhaus) y aplicaciones a series de Fourier o a familias de operadores.
    \item Teoremas de la aplicación abierta y de la gráfica cerrada. Equivalencia de normas que hacen de un espacio un Banach.
    \item Teorema de Hahn--Banach: versiones analítica y geométrica. Separación de convexos.
    \item Dual de $L^p$ y adjunto de un operador entre espacios de Banach.
\end{enumerate}

\item \textbf{Dualidad y topologías débiles (3 semanas):}
\begin{enumerate}
    \item Dual topológico, reflexividad y ejemplos ($L^p$, $\ell^p$, $c_0$).
    \item Topologías débil y débil-$*$. Convergencia débil frente a convergencia en norma.
    \item Teorema de Banach--Alaoglu. Compactidad débil y débil-$*$.
    \item Medidas de Radon y dual de $C(K)$ o de $C_0(X)$ (teorema de Riesz--Markov), al menos en el caso de $K$ compacto metrizable.
    \item Convexidad y teorema de Krein--Milman (\emph{ampliación}).
\end{enumerate}

\item \textbf{Operadores compactos y teorema espectral (2 semanas):}
\begin{enumerate}
    \item Operadores compactos y completamente continuos. Ideal de compactos. Ejemplos integrales.
    \item Alternativa de Fredholm en Hilbert (o en Banach, según el tiempo).
    \item Teorema espectral para operadores compactos autoadjuntos: descomposición espectral y convergencia de series de eigenfunciones.
    \item Aplicación a ecuaciones integrales de Fredholm o a problemas de Sturm--Liouville elementales.
\end{enumerate}

\end{enumerate}

\subsubsection*{Metodología}

\noindent \textbf{Lineamientos metodológicos:} La persona docente a cargo de este curso debe respetar las siguientes pautas:
\begin{enumerate}
    \item La presentación del material requiere de definiciones precisas, enunciados completos y demostraciones de los teoremas fundamentales (Hahn--Banach, acotación uniforme, aplicación abierta, gráfica cerrada, Alaoglu y teorema espectral compacto), aunque también es importante brindar contexto y motivación.
    \item Cada construcción abstracta debe acompañarse de al menos un espacio concreto ($L^p$, $\ell^p$, $C(K)$ o un espacio de Hilbert clásico) que ilustre el alcance de las hipótesis.
    \item Cada estudiante debe realizar una revisión de la bibliografía como complemento al trabajo presencial.
    \item Se fomenta la comunicación clara y rigurosa de argumentos, tanto oral como escrita.
\end{enumerate}

\textbf{Sugerencias metodológicas:} Adicionalmente, se tienen las siguientes recomendaciones para la persona docente a cargo de este curso:
\begin{enumerate}
    \item Explicitar el puente con el curso de análisis real: los espacios $L^p$ dejan de ser un apéndice de la integral de Lebesgue y pasan a ser el prototipo de espacio de Banach.
    \item Usar el espacio de Hilbert como laboratorio: proyección, Riesz y bases ortonormales deben aparecer antes de la teoría general de dualidad, para que Hahn--Banach se lea como su extensión.
    \item Dedicar tiempo a contraejemplos (operadores no acotados densamente definidos, normas no equivalentes, sucesiones débilmente convergentes que no convergen en norma) para evitar una lectura puramente formal de los teoremas.
    \item Promover el trabajo constante mediante listas de ejercicios que combinen demostraciones, cálculo de normas y construcción de ejemplos.
    \item Incorporar un proyecto o exposición breve sobre un tema de ampliación.
    \item Sugerir el uso de \LaTeX{} para la presentación de tareas y del proyecto.
\end{enumerate}


\nocite{Con90}
\nocite{Lax02}
\nocite{RS80}
\nocite{Bre11}
\nocite{Lan93}
\nocite{Rud91}
\nocite{Bog20}
\nocite{Fol99}

\printbibliography[heading=subbibliography,section=\the\value{refsection}]
\end{refsection}
